\section{Related Work}
\label{sec:related}

% Five subsections planned. Each cites the canonical papers.
%
% \subsection{Prompt compression for LLM context}
% LLMLingua and LongLLMLingua (Jiang et al., \cite{llmlingua,longllmlingua});
% RECOMP (\cite{recomp}); AutoCompressors (\cite{autocompressors}); ICAE
% (\cite{icae}). All evaluate on retrieval / QA tasks; none specifically on
% instruction documents. Compression-on-instructions is an open problem.
%
% \subsection{Hierarchical / progressive context}
% MemGPT (\cite{memgpt}) — three-tier memory (RAM / recall / archival);
% closest existing analog to our frequency-weighted residency model.
% Lost in the middle (Liu et al., \cite{lostinmiddle}) — established the
% U-shaped attention curve and the start-position bias for instruction-tuned
% models that motivates our "keep stable rules at top" recommendation.
%
% \subsection{Prompt caching and KV-cache reuse}
% Anthropic prompt-caching docs (\cite{anthropic-cache}); "Don't Break the
% Cache" (\cite{dont-break-cache}); PagedAttention (\cite{pagedattention});
% Mooncake (\cite{mooncake}); Preble (\cite{preble}). These motivate the
% caching baseline our methodology assumes.
%
% \subsection{Knowledge graphs for LLM context}
% GraphRAG (\cite{graphrag}); KAG (\cite{kag}). We argue these target
% relational corpora, not procedural instruction documents, and surface
% the serialisation overhead that rules them out for our setting.
%
% \subsection{Multi-agent coordination}
% AutoGen (\cite{autogen}); CrewAI; MetaGPT (\cite{metagpt}); ChatDev;
% LangGraph. These frame the problem domain; our contribution sits beneath
% the coordination layer at the per-agent context-loading level.

\tothink{Expand each subsection to $\sim$0.5 pages with proper citations.}
